\section{Exploratory development}
\label{sec:development}

This section reports the development phases: what the instrument was shown to be capable of
detecting, what an open adversary search found, and which candidate mechanisms were tested. All
measurements here are on \emph{training} material and are exploratory in the sense of
\S\ref{sec:evaluation-multiplicity}. None supports the paper's primary claim, and every pooled
interval reported in this section has a lower bound below the registered
\vsevenPrimaryFloor{}~pp reference bar, so each is a measured effect rather than a confirmed one.

\subsection{Baseline selection}
\label{sec:development-frontier}

\vsix{} deployed its policy with test-time search disabled, while also having measured that search
working. The configurations available at the end of that cycle therefore form what this paper
calls the \emph{\vsix{} frontier}: the set comprising the deployed non-searching policy and the
searching configurations \fcheap{} and \fmid{}, ordered by search cost. The term is the protocol's; it
denotes that comparison set of evaluated configurations, not a measured Pareto frontier in cost
and performance, and no dominance criterion was applied.

This distinction drove the most consequential methodological choice in the programme. Several
available configurations already exceeded the deployed policy, including unfitted single-switch
deviations of it, which motivated a searching configuration as the primary baseline. The
comparison target was fixed as \fcheap{}: the \vsix{} policy with its endgame-truncated search
enabled at the same search settings \sestina{} uses, and the least expensive searching member of
the frontier. \S\ref{sec:design-protocol} records that the choice was registered in advance.
Table~\ref{tab:configs} defines every configuration this paper compares.

\subsection{Detection calibration}
\label{sec:development-instrument}

The calibration procedure of \S\ref{sec:evaluation-floor} was built in this phase, which measured
the principal responder class's floor at \vsevenFloorPhaseOne{}~pp and predicted that it would fall
as (evaluation games)$^{-1/2}$. The phase of \S\ref{sec:development-adversaries} re-measured the
same handicap ladder at four times the games: the smallest recovered rung was then the
\vsevenPrimaryFloor{}~pp rung, and the \vsevenFloorRungBelow{}~pp rung remained unrecovered, so the
prediction was not borne out within the ladder's resolution. That is the first of the two findings
\S\ref{sec:evaluation-floor} reports, the second being that floors are family-specific, with the
declaration channel's at \vsevenDeclFloor{}~pp. Two further results shape later sections.

\paragraph{Transcript-inversion diagnostics.} Because the target
team is deterministic and published, an adversary can invert its transcript, asking for each observed
ask which deals would have produced it and contracting the posterior accordingly. Measured, this
contracts the deal posterior by approximately \vsevenBitsPerAsk{} bits per observed ask beyond the
certificates the rules already supply. Converting those bits into play is worth
\vsevenInversionUnfitted{}~pp unfitted against the deployed policy.

The gain was not target-specific. The excess tracks planted-edge size to within
$\pm$\vsevenInversionTrack{}~pp on all eight rungs of the calibration ladder, which is consistent
with a general strength gain rather than with a response that depends on a particular weakness of
the target; a broadly stronger policy can still be a valid profitable response, so this is a
statement about specificity rather than about exploitability. And no linear handicap adds more
than \vsevenHandicapBits{} bits per ask, so within this policy family no readability handicap was
found for a stochastic policy to remove. A proposed stochastic policy conditional on the inverter
succeeding was therefore not pursued.

\paragraph{Revised search-cost estimates.} The \vsix{} paper priced the searching configuration
at ``three orders of magnitude'' the blueprint, a ratio that mixed an all-threads timing with a
single-thread one. Re-timed here on one basis, it is \vsevenPhaseOneCost$\times$. But that figure
describes the \emph{unrestricted} search, which is not the frontier point anything was measured at,
and phase~3 re-measured the endgame-restricted operating point at approximately
\vsevenCostRecorded$\times$. \S\ref{sec:results-cost} corrects that correction in turn.

\subsection{Adversary generation against baseline configurations}
\label{sec:development-adversaries}

The constraint on this phase was that the searches must not share a bias, since repeated restarts
of one method do not provide methodological diversity. Every exploiter search in the corpus before it
had been the same search: a cross-entropy fit of a linear vector maximising win rate from one seed
against one target. This phase constructed six adversary classes in the synchronised regime, plus
one fitted search in the ex-ante-correlated regime, and, before fitting anything against them,
measured the effect of a diagnostic intervention on each channel under the evaluated condition: a
\vsix{} target, the phase-2 responder class and budget, and the training banks.

\paragraph{Adversarial-search outcome.}
Nothing constructed produced a resolved adversarial edge against the \vsix{} frontier beyond the
class floor. The largest adversarial response edge found against the deployed policy is an in-class
figure of \vsevenBestExploit{}~pp $[\vsevenBestExploitLo, \vsevenBestExploitHi]$ over
\vsevenBestExploitN{} games, which is below the class floor of \vsevenPrimaryFloor{}. Seven
separately fitted searches gave the same picture and none cleared its class floor.

Every configuration that beat the deployed policy by more than the floor proved on measurement not
to be target-specific: the target's own test-time search at \vsevenSearchAsExploit{}~pp, and a
single out-of-class coordinate at \vsevenOneSwitchArm{}~pp whose gain survives against opponents
sharing none of the target's machinery. The second of those is the contestation weight of
\S\ref{sec:agent-contest}; \S\ref{sec:results-attribution} measures on holdout what it contributes
and finds it the largest observed leave-one-out drop among the five components \sestina{} carries,
rather than the whole of the gain.

Two further results from this phase bear on \sestina{}'s design. The declaration-urgency channel
that the phase expected to be its headline has a diagnostic intervention estimate of
\vsevenUrgencyCeiling{}~pp $[\vsevenUrgencyCeilingLo, \vsevenUrgencyCeilingHi]$ under that
condition and did not yield a resolved adversarial advantage at the tested budget, but the same channel is worth \vsevenUrgencyWorthLo{} to
\vsevenUrgencyWorthHi{}~pp to a configuration that simply disables it: the adversary owns the
smaller half of a two-point defect in the target. And the event-count cliff of
\S\ref{sec:agent-urgency} was not reached by any adversary in the tested searches, while a
configuration that removes the counter inherits a self-play game-length tail of \vsevenCliffTail{}
events, which is the obligation \S\ref{sec:agent-stall} discharges.

\subsection{Candidate-mechanism evaluation}
\label{sec:development-rejected}

Candidate modifications were evaluated on development data before configuration freezing. Five
were developed in isolated branches, gated before being scored, and checked against the
side-channel definition by the harness of \S\ref{sec:evaluation-side} rather than by inspection.
One was retained, and it is the configuration change described in \S\ref{sec:agent} rather than an
architecture. Four did not demonstrate an additional performance benefit under the tested settings:
alternative declaration allocation, learned leaf evaluation, per-decision training objectives, and
learned action ranking. The following subsections report their measured outcomes and limitations.
These results concern the tested implementations and budgets, not all methods in the corresponding
classes.

\subsubsection{Declaration-allocation comparison}
\label{sec:development-alloc}

The design register's second-ranked hypothesis, at priority \vsevenLonePriority{}, held that
\vsix{}'s entire margin lives in the declaration channel and that roughly two points of it remained.
Its proposed mechanism was that the deployed allocator, a product of per-card marginals, selects a
different allocation from the joint posterior's argmax often enough to matter.

The measurement did not support that mechanism. Over \vsevenLoneDeclsTotal{} recorded voluntary
declarations across two banks and both urgency settings, \vsevenLoneDeclsDistinct{} of them from
distinct (bank, urgency-setting) capture cells and the remainder from re-runs of the same seeds and
from the two cells of the sequential-joint policy, the sequential approximate joint's argmax and
the marginal-product argmax selected the same allocation at every one, with the joint score holding a
strict opinion on \vsevenLoneStrictLo\% to \vsevenLoneStrictHi\% of ambiguous cases.
This agreement is empirical: the two procedures are not equivalent in general, since the six cards
of a half-suit are coupled through the capacity constraint, and no argument that the sequential
conditioning preserves the marginal-product ranking is supplied.

Two further measurements bear on the direction rather than only the mechanism. First, the exact
constraint-joint maximiser \eqref{eq:bestalloc} performed worse than the deployed approximation on
the declaration decision itself: replayed on one capture cell of \vsevenJointCellDecls{} voluntary
declarations, it corrects \vsevenJointFixed{} declarations and breaks \vsevenJointBroken{}, a net
\vsevenJointWorseNet{}~pp of declaration accuracy on that cell. At the design register's fitted
conversion of about \vsevenDeclToWinConv{} win-rate points per point of declaration accuracy
(range \vsevenDeclToWinConvLo{} to \vsevenDeclToWinConvHi{} across banks), this corresponds to
approximately \vsevenJointWorseWin{}~pp of win rate; that figure is inferred from the
declaration-level difference, not measured as a game outcome. This is the second comparison of the
kind \S\ref{sec:background-vsix} reports. Second, \vsevenJointFlatLo\% to \vsevenJointFlatHi\% of
the error class occupies states in which the exact constraint-conditioned distribution assigns the
competing allocations equal probability, so most of the apparent headroom is an information limit
under that model rather than a policy defect.

The channel is also largely downstream of a defect that \sestina{} removes for other reasons.
Disabling the urgency escalation eliminates \vsevenLoneUrgencyShare\% of the entire
allocation-error class and drops the oracle-assisted estimate for a perfect allocator, at the
measured incidence, from \vsevenLoneCeilingBefore{} to \vsevenLoneCeilingAfter{}~pp, which is below
the reference bar. The
consequence is a composition constraint this paper observes: a declaration-channel gain and the
urgency-off gain should not be added, because they overlap in the same error class.

\subsubsection{Learned leaf evaluation}
\label{sec:development-leaf}

An inherited hypothesis held that the test-time search cannot be judged until its leaf evaluator is
rebuilt, the deployed one being algebraically close to a rescaling of the hit probability. A leaf was
fitted for the purpose, and the statistic that matters, between-candidate $R^2$, which is the only
component the paired deviation rule can consume, improved from \vsevenKoneRsqBefore{} to
\vsevenKoneRsqAfter{} in sample and from \vsevenKoneRsqOosBefore{} to \vsevenKoneRsqOosAfter{} out
of sample; the in-sample ratio is \vsevenKoneFold-fold, the out-of-sample ratio is much smaller,
and the absolute $R^2$ remains modest in both, so the direction replicates and the multiple does
not. The
fit was scored over approximately \vsevenKoneLeaves{} leaves across two search regimes and three
banks.

In play it is worth \vsevenKoneInPlay{}~pp against the material control's \vsevenKoneControl{}~pp,
losing on both banks individually. The diagnosis is conversion failure rather than fit failure: the
rule's deviation rate is \vsevenKoneChangeLo{} for the deployed material leaf,
\vsevenKoneChangeMid{} for the fitted one and \vsevenKoneChangeHi{} for a third arm that merely
rescales the material leaf's constant, so at the deployed shrinkage the deviation rule is nearly
leaf-invariant. The tested leaf replacement did not improve playing strength under the deployed
selection rule. This illustrates a mismatch between a surrogate metric and playing performance
related to the abstraction pathology the poker literature documents \cite{waugh-abstraction}, in
which a model improved on its own metric can leave play unchanged or worse; no formal mapping
between the two settings is claimed.

The inherited hypothesis was discharged separately and by re-measurement rather than by the new
evaluator. The single underpowered cell on which it rested, one \vsevenFullGameOldN{}-game cell at
\vsevenFullGameOld{}~pp, is the \texttt{depth=24} configuration; re-run on the training banks with
\vsix{}'s own leaf it gives \vsevenDepthTwentyFour{}
$[\vsevenDepthTwentyFourLo, \vsevenDepthTwentyFourHi]$, and full-game search at \texttt{depth=12}
gives \vsevenFullGameReOne{} and \vsevenFullGameReTwo{}~pp, pooled \vsevenFullGameRePooled{}.
Phase 1's \vsevenTruncatedEndgame{}~pp refuted the conditional in the endgame regime only, at
\texttt{depth=12} with \texttt{maxq=26}; it was this full-game cell that left it binding for
full-game search.

\subsubsection{Per-decision training objectives}
\label{sec:development-perdecision}

A further hypothesis held that a per-decision objective is necessary because the scoreboard is too
noisy to resolve small effects. Its precision claim is supported and its fitting form is not, and
the two halves are independent.

The design effect of deal clustering was measured for the first time in this cycle. A
one-point-equivalent effect reaches a planning half-width of its own size in
\vsevenPerDecisionGamesLo{} to \vsevenPerDecisionGamesHi{} games on the declaration channel and
\vsevenAllocGamesLo{} to \vsevenAllocGamesHi{} on the allocation-error channel, against
\vsevenScoreboardGames{} on the scoreboard, a ratio of \vsevenDesignEffectLo{} to
\vsevenDesignEffectHi{} times. That ratio compares planning sample sizes for different outcome
channels, linked by the design register's assumed conversion factors between per-decision accuracy
and win rate; the per-decision figures apply the measured clustering design effect and the
scoreboard figure is the independent-game approximation of \S\ref{sec:evaluation-power}, so the
ratio is a planning comparison rather than a measured precision gain on one estimand.

But all three self-oriented per-decision objectives, fitted at matched budget against a per-game
control, produced worse policies: \vsevenKfourSelfDecl{}, \vsevenKfourSelfAsk{} and
\vsevenKfourSelfAlloc{}~pp, each replicated in sign and two of them with intervals beyond the
reference bar. Each fit moved its own proxy in the intended direction and lost anyway. The
declaration-oriented fit increased declaration accuracy by \vsevenKfourTradeGain{}~pp while
reducing ask accuracy by \vsevenKfourTradeCost{}~pp, and the ask-oriented fit did the reverse. The
two per-decision channels exhibited a trade-off in these fits, so a fitter holding
\vsevenCoordinates{} coordinates improved one at the expense of the other. The results motivate
further study of the estimator without establishing the tested proxy as a suitable training
objective.

\subsubsection{Learned action ranking and tie-breaking}
\label{sec:development-tie}

In the \vsix{} mirror \vsevenTieContested\% of contested ask decisions tie bit-for-bit, and the
open question was whether anything can be \emph{selected} within that set rather than merely
randomised. The learned candidate this cycle built was not built for that question. It was the
brief's wildcard, aimed at amortising test-time search, and it was fitted unrestricted over
\vsevenKfiveFitDecisions{} searched decisions. Deployed unrestricted it scores
\vsevenKfiveUnrestricted{}~pp pooled over \vsevenPhaseThreeCell{} games, negative on both banks
$[\vsevenKfiveUnrestrictedLo, \vsevenKfiveUnrestrictedHi]$. Restricting it to the tied set is one
added switch on the same fitted weight vector, and it is there that the tie-group question gets
answered.

Restricted, against the deployed policy it scores \vsevenKfiveTieOnly{}~pp, not statistically
resolved from the free public-history tie-break's \vsevenKfiveFreeTie{}~pp; head to head against
that free tie-break at \vsevenPhaseThreeCell{} games it is \vsevenKfiveHeadToHead{}~pp $\pm$
\vsevenKfivePm{} (95\% half-width), replicated in sign. The restricted learned ranker did not
demonstrate an advantage over public-history tie-breaking, and the appropriate control for any
mechanism operating inside the tied set is the decorrelated policy rather than the deployed one.

The same instrumentation explains why. Setting every candidate's true advantage to exactly zero and
feeding the capture's real standard errors through the engine's real deviation rule reproduces the
search's observed deviation rate to a \vsevenCurseRelErr\% relative error, on a quantity the null was
never fitted to. The search's per-decision selection signal is approximately \vsevenCurseShare\%
winner's curse on its own Monte Carlo draw, and approximately \vsevenCurseTieShare\% within the tied
set, where the rule applies no shrinkage at all.

\subsubsection{Forced-endgame improvement estimates}
\label{sec:development-arithmetic}

Closing the entire forced-endgame accuracy gap at the measured incidence, re-derived from this
cycle's own capture, is estimated at \vsevenLthirteenClose{}~pp, which is the register's own figure
corrected upward by a factor of two and still a fiftieth of the reference bar. That estimate holds
for the incidence observed in the measured play. The route by which the estimate could grow is an
adversary driving the target's incidence upward, and the one construction measured does not do
that: a team that never declares voluntarily raises its \emph{own} forced-endgame rate, not its
opponent's, and pays \vsevenLthirteenAdvCostLo{} to \vsevenLthirteenAdvCostHi{}~pp for it. No
evaluated intervention increased the target's forced-endgame incidence, so this route remains
untested rather than closed (Appendix~\ref{app:corrections-self} records the earlier report's
unsupported figure for it).

\subsection{Configuration selection and freezing}
\label{sec:development-freeze}

\vsevenSelK{} configurations were scored against the deployed policy during development. One was
selected, its specification written to an artifact whose reconstruction is executed rather than
asserted, and the configuration frozen. Nothing about it was changed afterwards, and no gradient or population fit was applied to it. \S\ref{sec:results-selection} quantifies what that selection could be worth
under the null.
